% Mathematical fragments are supplied by derivation_math.tex.
% Condensed 2026-10-31: purely intermediate algebra lines were removed; every
% numbered equation, its number and its order are retained.
\section{复 Hamiltonian 表示与辛差分格式构造}
\label{sec:fixed_method}

承接有限维投影系统，本节先由实Hamiltonian导出复坐标下的单生成量表示并界定其对称与有势条件，再针对构型依赖质量导致的不可分离时间推进问题，回到实正则空间借助Birkhoff生成函数与中点变换构造隐式更新式（36）。

\subsection{有限维系统的 Hamiltonian 表示}

在对称可逆质量矩阵$M(q)$下，由广义动量
\DU{P0149M01}
得动能$T=\tfrac12p^TM^{-1}(q)p$；记伸长与弯曲势能之和为$V=U_{\mathrm{stretch}}+U_{\mathrm{bend}}$（梯度为$K(q)q$），系统实能量为
\DU{P0169M01}
将未并入$V$的外载荷与流体阻力作为非保守广义力引入动量方程，状态演化与功率收支依次写为
\DU{P0160M01}
\DU{P0165M01}
\DE{P0170M01}{12}
式（12）将能量变化归结为剩余外力做功与阻尼耗散。变质量时$\partial H/\partial q$含$M^{-1}(q)$的构型导数；格式诊断保留该项，两模态梁学习取常质量分支。

\subsection{复坐标与复 Hamiltonian}

为将含耗散与外力的演化收归单一复函数导数，定义复生成量
\DE{P0337M01}{16}
其中$H$为未含剩余外力势的实能量，$G_1$为耗散生成函数，$G_2$为剩余外力势（已入$H$的静态磁势不重复计入）。常系数线性阻尼下设$M=M^T\succ0$、$C=C^T\succeq0$且$A=CM^{-1}=A^T$，使阻尼力表为$G_1$梯度：
\DU{P0341M01}
忽略非光滑接触且剩余外力有标量势时，取
\DU{P0344M01}

引入对称归一化的复坐标与Wirtinger导数：
\DE{P0349M01}{17}
\DE{P0350M01}{18}
由逆变换知辛二形式为$\omega=dq\wedge dp=i\,d\xi\wedge d\bar\xi$。原状态方程及保守梯度相应化为
\DE{P0366M01}{19}
\DE{P0367M01}{20}
代入式（18）合并$G_1$、$G_2$贡献，得
\DE{P0423M01}{21}
其中$\mathcal P$简记方括号内载荷项。再引入辅助势
\DU{P0576M01}
即可并入同一共轭导数：
\DE{P0571M01}{25}
此处$\widetilde H=\mathcal H$即式（16）。式（25）表明，在对称阻尼与有势外力条件下，复配对$\xi=(q+ip)/\sqrt2$将状态演化紧凑写为复生成量的共轭导数，这为第\ref{sec:experiments}节复值Hamiltonian网络提供了坐标依据（后续学习取保守分支$\mathcal H=H+i0$并在实正则坐标中求梯度；一般构型或速度相关耗散不直接适用式（25））。然而，在推进轨迹与生成数据时，若$M(q)$依赖构型导致动能不可分离，通常按可分离动能与势能构造的显式蛙跳更新不能直接沿用；因此下文回到实正则空间，借Birkhoff生成函数理论构造保守主部的保辛隐式中点格式，再并入独立中点非保守载荷。

\subsection{Birkhoff 表示与生成函数}

针对$H$的保守主部，采用Birkhoff系统框架，其运动方程为
\DU{P0796M01}
其中$B$为Birkhoff函数，$F_i$为结构函数。按$(q,p)$分块取$F_q=p$、$F_p=0$（一自由度即$F_1=p$、$F_2=0$），保守主部对应于
\DU{P0850M01}

在系统正则、自伴且变换相容条件下，以双点变换
\DU{P0854M01}
联系固定初态$z$与当前态$\hat z$。记其Jacobian上下块行为$(A_\alpha,B_\alpha)$与$(C_\alpha,D_\alpha)$；若满足局部非退化条件
\DU{P0866M01}
则小步长下$\hat w=f(w,t,t_0)$具对称Jacobian，可由生成函数梯度表示：
\DU{P0870M01}
其偏微分演化关系为
\DU{P0871M01}
其中$A$由原变量算子$\overline A$经逆变换得到，而$\overline A$直接联系变换变量导数：
\DU{P0873M01}
当$A$与$\alpha$解析时，生成函数梯度具收敛级数展开：
\DU{P0884M01}
零阶项由恒等初值确定，首阶系数按
\DU{P0886M01}
计算。上述生成函数存在性定理\citep[第15章，定理4.1]{FengQin2010}、幂级数递推定理\citep[第15章，定理4.4]{FengQin2010}及生成函数梯度的$m$阶截断定理\citep[第15章，定理4.5]{FengQin2010}，其陈述与证明均归属所引专著。在前述正则、变换相容与局部非退化条件下，梯度截断可定义保守主部的辛格式，无需动能与势能可分离；截断所得格式是否具有时间对称性，仍须由具体构造判断。下文选定中点变换求取保守首阶系数，并单独加载非保守力。

\subsection{中点变换与生成函数系数}

固定初刻$t_0$及初态$z=(q,p)$，以$\delta=t-t_0$参数化末态$\hat z=(\hat q,\hat p)$，取中点$w=(Q,P)$与增量$\hat w=(\hat Q,\hat P)$，变换、Jacobian及逆变换依次为
\DE{P0909M01}{30}
\DU{P0916M01}
\DE{P0919M01}{31}
因$C_\alpha+D_\alpha=\operatorname{diag}(I,-I)$非奇异，局部非退化条件恒成立。固定初态对当前时刻求导，得
\DE{P0930M01}{32}
\DU{P0931M01}
在$\delta=0$时，$\hat z=z$使$\varphi_w^{(0)}=0$且$\partial\hat w/\partial w=0$，故$\varphi_w^{(1)}$等于初刻$d\hat w/dt$的保守部分；该系数与同中点独立载荷分别为
\DU{P0959M01}
将保守生成函数梯度截断至首阶，再在同一中点叠加独立载荷，定义一步更新：
\DE{P0960M01}{33}
式（33）分离了保守生成梯度与非保守载荷，后者不沿用保守流的保辛保证。保守部分取$m=1$时，定理4.5至少保证一阶精度；具体实例的更高观察阶由第\ref{sec:experiments}节的步长实验检验。

\subsection{返回状态变量与差分更新}

以$(q_k,p_k)$和$(q_{k+1},p_{k+1})$替代$(q,p)$与$(\hat q,\hat p)$，返回原正则变量：
\DE{P0962M01}{34}
其中中点速度取$v=-M^{-1}(Q)P$。略去中点自变量并以$\Delta$记变质量收缩项，得
\DE{P0985M01}{35}
由首式消去位置增量并引入变质量二次收缩与矩阵$\epsilon$：
\DU{P0994M01}
经矩阵左整理即得最终更新式
\DE{P1002M01}{36}
式（36）中$(I+\epsilon)^{-1}$从左侧作用，各系数在位置与动量中点取值：变系数时隐式求解，常质量时$\Delta=0$。保守分支依前述定理保辛，含非保守载荷时的全映射结构则在第\ref{sec:experiments}节另行核验。
